\section{Noetherian graded rings}
\label{section-noetherian-graded}

\noindent
A bit of theory on Noetherian graded rings including some material on
Hilbert polynomials.

\begin{lemma}
\label{lemma-S-plus-generated}
Let $S$ be a graded ring. A set of homogeneous elements
$f_i \in S_{+}$ generates $S$ as an algebra over $S_0$ if
and only if they generate $S_{+}$ as an ideal of $S$.
\end{lemma}

\begin{proof}
If the $f_i$ generate $S$ as an algebra over $S_0$ then every element
in $S_{+}$ is a polynomial without constant term in the $f_i$ and hence
$S_{+}$ is generated by the $f_i$ as an ideal. Conversely, suppose that
$S_{+} = \sum Sf_i$. We will prove that any element $f$ of $S$ can be written
as a polynomial in the $f_i$ with coefficients in $S_0$. It suffices
to do this for homogeneous elements. Say $f$ has degree $d$. Then we may
perform induction on $d$. The case $d = 0$ is immediate. If $d > 0$
then $f \in S_{+}$ hence we can write $f = \sum g_i f_i$
for some $g_i \in S$. As $S$ is graded we can replace $g_i$ by its
homogeneous component of degree $d - \deg(f_i)$. By induction we
see that each $g_i$ is a polynomial in the $f_i$ and we win.
\end{proof}

\begin{lemma}
\label{lemma-graded-Noetherian}
A graded ring $S$ is Noetherian if and only if $S_0$ is
Noetherian and $S_{+}$ is finitely generated as an ideal of $S$.
\end{lemma}

\begin{proof}
It is clear that if $S$ is Noetherian then $S_0 = S/S_{+}$ is Noetherian
and $S_{+}$ is finitely generated. Conversely, assume $S_0$ is Noetherian
and $S_{+}$ finitely generated as an ideal of $S$. Pick generators
$S_{+} = (f_1, \ldots, f_n)$. By decomposing the $f_i$ into homogeneous
pieces we may assume each $f_i$ is homogeneous. By
Lemma \ref{lemma-S-plus-generated}
we see that $S_0[X_1, \ldots X_n] \to S$ sending $X_i$ to $f_i$
is surjective. Thus $S$ is Noetherian by
Lemma \ref{lemma-Noetherian-permanence}.
\end{proof}

\begin{definition}
\label{definition-numerical-polynomial}
Let $A$ be an abelian group.
We say that a function $f : n \mapsto f(n) \in A$
defined for all sufficient large integers $n$ is a
{\it numerical polynomial} if there exists $r \geq 0$,
elements $a_0, \ldots, a_r\in A$ such that
$$
f(n) = \sum\nolimits_{i = 0}^r \binom{n}{i} a_i
$$
for all $n \gg 0$.
\end{definition}

\noindent
The reason for using the binomial coefficients is the
elementary fact that any polynomial $P \in \mathbf{Q}[T]$
all of whose values at integer points are integers, is
equal to a sum $P(T) = \sum a_i \binom{T}{i}$ with
$a_i \in \mathbf{Z}$. Note that in particular the
expressions $\binom{T + 1}{i + 1}$ are of this form.

\begin{lemma}
\label{lemma-numerical-polynomial-functorial}
If $A \to A'$ is a homomorphism of abelian groups and if
$f : n \mapsto f(n) \in A$ is a numerical polynomial,
then so is the composition.
\end{lemma}

\begin{proof}
This is immediate from the definitions.
\end{proof}

\begin{lemma}
\label{lemma-numerical-polynomial}
Suppose that $f: n \mapsto f(n) \in A$
is defined for all $n$ sufficiently large
and suppose that $n \mapsto f(n) - f(n-1)$
is a numerical polynomial. Then $f$ is a
numerical polynomial.
\end{lemma}

\begin{proof}
Let $f(n) - f(n-1) = \sum\nolimits_{i = 0}^r \binom{n}{i} a_i$
for all $n \gg 0$. Set
$g(n) = f(n) - \sum\nolimits_{i = 0}^r \binom{n + 1}{i + 1} a_i$.
Then $g(n) - g(n-1) = 0$ for all $n \gg 0$. Hence $g$ is
eventually constant, say equal to $a_{-1}$. We leave it
to the reader to show that
$a_{-1} + \sum\nolimits_{i = 0}^r \binom{n + 1}{i + 1} a_i$
has the required shape (see remark above the lemma).
\end{proof}

\begin{lemma}
\label{lemma-graded-module-fg}
If $M$ is a finitely generated graded $S$-module,
and if $S$ is finitely generated over $S_0$, then
each $M_n$ is a finite $S_0$-module.
\end{lemma}

\begin{proof}
Suppose the generators of $M$ are $m_i$ and the generators
of $S$ are $f_i$. By taking homogeneous components we may
assume that the $m_i$ and the $f_i$ are homogeneous
and we may assume $f_i \in S_{+}$. In this case it is
clear that each $M_n$ is generated over $S_0$
by the ``monomials'' $\prod f_i^{e_i} m_j$ whose
degree is $n$.
\end{proof}

\begin{proposition}
\label{proposition-graded-hilbert-polynomial}
Suppose that $S$ is a Noetherian graded ring
and $M$ a finite graded $S$-module. Consider the
function
$$
\mathbf{Z} \longrightarrow K'_0(S_0), \quad
n \longmapsto [M_n]
$$
see Lemma \ref{lemma-graded-module-fg}.
If $S_{+}$ is generated by elements of degree $1$,
then this function is a numerical polynomial.
\end{proposition}

\begin{proof}
We prove this by induction on the minimal number of
generators of $S_1$. If this number is $0$, then
$M_n = 0$ for all $n \gg 0$ and the result holds.
To prove the induction step, let $x\in S_1$
be one of a minimal set of generators, such that
the induction hypothesis applies to the
graded ring $S/(x)$.

\medskip\noindent
First we show the result holds if $x$ is nilpotent on $M$.
This we do by induction on the minimal integer $r$ such that
$x^r M  = 0$. If $r = 1$, then $M$ is a module over $S/xS$
and the result holds (by the other induction hypothesis).
If $r > 1$, then we can find a short exact sequence
$0 \to M' \to M \to M'' \to 0$ such that the integers
$r', r''$ are strictly smaller than $r$. Thus we know
the result for $M''$ and $M'$. Hence
we get the result for $M$ because of the relation
$
[M_d]  = [M'_d] + [M''_d]
$
in $K'_0(S_0)$.

\medskip\noindent
If $x$ is not nilpotent on $M$, let $M' \subset M$ be
the largest submodule on which $x$ is nilpotent.
Consider the exact sequence $0 \to M' \to M \to M/M' \to 0$
we see again it suffices to prove the result for $M/M'$. In other
words we may assume that multiplication by $x$ is injective.

\medskip\noindent
Let $\overline{M} = M/xM$. Note that the map $x : M \to M$
is {\it not} a map of graded $S$-modules, since it does
not map $M_d$ into $M_d$. Namely, for each $d$ we have the
following short exact sequence
$$
0 \to M_d \xrightarrow{x} M_{d + 1} \to \overline{M}_{d + 1} \to 0
$$
This proves that $[M_{d + 1}] - [M_d] = [\overline{M}_{d + 1}]$.
Hence we win by Lemma \ref{lemma-numerical-polynomial}.
\end{proof}

\begin{remark}
\label{remark-period-polynomial}
If $S$ is still Noetherian but $S$ is not generated in degree $1$,
then the function associated to a graded $S$-module is a periodic
polynomial (i.e., it is a numerical polynomial on the
congruence classes of integers modulo $n$ for some $n$).
\end{remark}

\begin{example}
\label{example-hilbert-function}
Suppose that $S = k[X_1, \ldots, X_d]$.
By Example \ref{example-K0-field} we may identify
$K_0(k) = K'_0(k) = \mathbf{Z}$. Hence any finitely
generated graded $k[X_1, \ldots, X_d]$-module
gives rise to a numerical polynomial
$n \mapsto \dim_k(M_n)$.
\end{example}

\begin{lemma}
\label{lemma-quotient-smaller-d}
Let $k$ be a field. Suppose that $I \subset k[X_1, \ldots, X_d]$
is a nonzero graded ideal. Let $M = k[X_1, \ldots, X_d]/I$.
Then the numerical polynomial $n \mapsto \dim_k(M_n)$ (see
Example \ref{example-hilbert-function})
has degree $ < d - 1$ (or is zero if $d = 1$).
\end{lemma}

\begin{proof}
The numerical polynomial associated to the graded module
$k[X_1, \ldots, X_d]$ is $n \mapsto \binom{n - 1 + d}{d - 1}$.
For any nonzero homogeneous $f \in I$ of degree $e$
and any degree $n >> e$ we have $I_n \supset f \cdot k[X_1, \ldots, X_d]_{n-e}$
and hence $\dim_k(I_n) \geq \binom{n - e - 1 + d}{d - 1}$. Hence
$\dim_k(M_n) \leq \binom{n - 1 + d}{d - 1} - \binom{n - e - 1 + d}{d - 1}$.
We win because the last expression
has degree $ < d - 1$ (or is zero if $d = 1$).
\end{proof}








